\begin{textsample}{POS Dim 3 – rj – Score 13.00 – rj/rj\_inf\_cam\_exp\_11.txt}  \label{ex:f3_pos_018}
A capacidade de \textbf{criança} construir representações e expressá -las pela linguagem verbal se desenvolve conforme a \textbf{criança} interage com pessoas mais experientes e busca captar os signos e símbolos construídos socialmente, presentes nos comportamentos de seus parceiros humanos e expressos na oralidade ou na escrita no ambiente em que ela convive.
A linguagem oral ( ou a língua de sinais ) atravessa todo o cotidiano das \textbf{crianças} na Educação \textbf{Infantil}.
Através da oralidade ( ou da vocalização e sinalização, no caso de \textbf{crianças} com deficiência auditiva ) as \textbf{crianças} expressam vontades, \textbf{desejos}, fazem perguntas, \textbf{contam} casos, concordam ou discordam de um \textbf{colega} ou da \textbf{professora} e são muito interessadas no efeito que suas manifestações verbais provocam em outras pessoas.
Também as falas dos parceiros \textbf{adultos} e \textbf{crianças}, de artistas em programas de televisão ou de personagens das histórias lidas são objeto de observação curiosa e de \textbf{imitação} pelas \textbf{crianças} e lhes abrem um universo maior de fontes de apropriação da oralidade ( ou língua de sinais ).
Sua aquisição assegura às \textbf{crianças} a possibilidade de participar de situações cotidianas onde podem comunicar -se, \textbf{conversar}, ouvir histórias, narrar, \textbf{contar} um fato, \textbf{brincar} com palavras, expressar sua opinião, comparar conceitos, construindo estratégias para conhecer o mundo.
A \textbf{criança}, por sua vez, conforme amplia suas experiências na cultura, percebe que muitas vezes os \textbf{adultos} com quem convive se ocupam de compreender instruções em \textbf{embalagens}, fazer pedidos ou dar informações por meio de grafismos, que se referem à linguagem escrita, carregada de novas características, que estimulam a \textbf{criança} a buscar compreender seu funcionamento no contexto em que vive.
A denominação deste campo busca evidenciar a estreita relação entre os atos de falar e escutar com a constituição da linguagem e do pensamento humanos, desde a infância.
A aproximação de diferentes linguagens traz para o cotidiano das unidades de Educação \textbf{Infantil} momentos de “ escutar ”, no sentido de produzir/acolher mensagens orais, gestuais, corporais, musicais, \textbf{plásticas}, além das mensagens trazidas por textos escritos, e “ falar ”, entendido como expressar/interpretar não apenas pela oralidade, mas também pela linguagem de sinais, pela escrita convencional ou não convencional, pela escrita braile, e também pelas danças, desenhos e outras manifestações expressivas.
Escutar e falar não se restringem a um só campo de experiências, mas são atos transversais a todos os campos, embora aqui apresentado com mais profundidade.
No sentido de dar força ao olhar para a pluralidade de linguagens que preside a educação em geral, e a Educação \textbf{Infantil} em particular, e considerando que a linguagem verbal não se separa completamente da linguagem corporal, musical, \textbf{plástica} e \textbf{dramática}, será feita a necessária aproximação entre o campo escuta, fala, pensamento e imaginação e os demais campos de experiências.
Para tanto \textbf{contamos} com o apoio das DCNEI ( Parecer CNE/CEB nº 20/09 ), para quem as práticas pedagógicas vividas na Educação \textbf{Infantil} devem garantir experiências que"favoreçam a imersão das \textbf{crianças} nas diferentes linguagens e o progressivo domínio por elas de vários gêneros e formas de expressão : gestual, verbal, \textbf{plástica}, \textbf{dramática} e musical".

% matched lemmas: adulto, brincar, colega, contar, conversar, criança, desejo, dramático, embalagem, imitação, infantil, plástico, professora
\end{textsample}
